\documentclass[11pt]{article}
\usepackage[margin=1in]{geometry}
\usepackage{graphicx}
\usepackage{booktabs}
\usepackage{amsmath,amssymb}
\usepackage{xcolor}
\usepackage{caption}
\usepackage[colorlinks=true,linkcolor=blue!50!black,urlcolor=blue!50!black]{hyperref}
\captionsetup[table]{name=Supplementary Table,labelsep=colon}
\renewcommand{\thetable}{S\arabic{table}}
\newcommand{\SubOnly}{0.580}
\newcommand{\SubOnlyLo}{0.565}
\newcommand{\SubOnlyHi}{0.594}
\newcommand{\TabNPath}{22,254}
\newcommand{\TabNCommon}{17,396}
\newcommand{\PureNPath}{16,903}
\newcommand{\PureNCommon}{1,479}
\newcommand{\PurePctPath}{76.0}
\newcommand{\PurePctCommon}{8.5}
\newcommand{\SampPurePctPath}{74.9}
\newcommand{\SampPurePctCommon}{7.9}
\newcommand{\NSample}{6,000}
\newcommand{\NPerClass}{3,000}
\newcommand{\MetaAll}{0.968}
\newcommand{\MetaAllLo}{0.963}
\newcommand{\MetaAllHi}{0.972}
\newcommand{\MetaHarm}{0.835}
\newcommand{\MetaHarmLo}{0.824}
\newcommand{\MetaHarmHi}{0.845}
\newcommand{\MetaFlag}{0.812}
\newcommand{\MetaFlagLo}{0.801}
\newcommand{\MetaFlagHi}{0.823}
\newcommand{\MetaTerms}{0.958}
\newcommand{\MetaTermsLo}{0.953}
\newcommand{\MetaTermsHi}{0.964}
\newcommand{\MetaShared}{0.624}
\newcommand{\MetaSharedLo}{0.609}
\newcommand{\MetaSharedHi}{0.637}
\newcommand{\DnbRef}{0.681}
\newcommand{\DnbRefLo}{0.668}
\newcommand{\DnbRefHi}{0.694}
\newcommand{\DnbDelta}{0.507}
\newcommand{\DnbDeltaLo}{0.492}
\newcommand{\DnbDeltaHi}{0.522}
\newcommand{\HynRef}{0.617}
\newcommand{\HynRefLo}{0.604}
\newcommand{\HynRefHi}{0.632}
\newcommand{\HynDelta}{0.524}
\newcommand{\HynDeltaLo}{0.508}
\newcommand{\HynDeltaHi}{0.537}
\newcommand{\KmRef}{0.568}
\newcommand{\KmRefLo}{0.554}
\newcommand{\KmRefHi}{0.583}
\newcommand{\KmDelta}{0.504}
\newcommand{\KmDeltaLo}{0.490}
\newcommand{\KmDeltaHi}{0.519}
\newcommand{\PureStratP}{2,247}
\newcommand{\PureStratC}{236}
\newcommand{\DnbRefPure}{0.749}
\newcommand{\DnbRefPureLo}{0.711}
\newcommand{\DnbRefPureHi}{0.782}
\newcommand{\DnbDeltaPure}{0.493}
\newcommand{\DnbDeltaPureLo}{0.455}
\newcommand{\DnbDeltaPureHi}{0.532}
\newcommand{\HynRefPure}{0.635}
\newcommand{\HynRefPureLo}{0.596}
\newcommand{\HynRefPureHi}{0.674}
\newcommand{\HynDeltaPure}{0.527}
\newcommand{\HynDeltaPureLo}{0.487}
\newcommand{\HynDeltaPureHi}{0.563}
\newcommand{\KmRefPure}{0.552}
\newcommand{\KmRefPureLo}{0.512}
\newcommand{\KmRefPureHi}{0.591}
\newcommand{\KmDeltaPure}{0.460}
\newcommand{\KmDeltaPureLo}{0.422}
\newcommand{\KmDeltaPureHi}{0.501}
\newcommand{\MetaPure}{0.582}
\newcommand{\MetaPureLo}{0.545}
\newcommand{\MetaPureHi}{0.620}
\newcommand{\DnbRefPureTrained}{0.721}
\newcommand{\DnbRefPureTrainedLo}{0.683}
\newcommand{\DnbRefPureTrainedHi}{0.757}
\newcommand{\HynRefPureTrained}{0.649}
\newcommand{\HynRefPureTrainedLo}{0.609}
\newcommand{\HynRefPureTrainedHi}{0.688}
\newcommand{\MetaNC}{1.000}
\newcommand{\MetaNCLo}{1.000}
\newcommand{\MetaNCHi}{1.000}
\newcommand{\DnbRefNC}{0.633}
\newcommand{\DnbRefNCLo}{0.607}
\newcommand{\DnbRefNCHi}{0.662}
\newcommand{\NCP}{424}
\newcommand{\NCC}{1,897}
\newcommand{\IncDnb}{0.010}
\newcommand{\IncDnbLo}{0.007}
\newcommand{\IncDnbHi}{0.013}
\newcommand{\IncHyn}{0.005}
\newcommand{\IncHynLo}{0.002}
\newcommand{\IncHynHi}{0.007}
\newcommand{\ShufMean}{0.501}
\newcommand{\ShufMax}{0.519}
\newcommand{\AgMeta}{0.959}
\newcommand{\AgMetaLo}{0.945}
\newcommand{\AgMetaHi}{0.970}
\newcommand{\AgRef}{0.608}
\newcommand{\AgRefLo}{0.573}
\newcommand{\AgRefHi}{0.643}
\newcommand{\AgDelta}{0.514}
\newcommand{\AgDeltaLo}{0.481}
\newcommand{\AgDeltaHi}{0.552}
\newcommand{\AgN}{1,018}
\newcommand{\AgPath}{461}
\newcommand{\AgCommon}{557}
\newcommand{\CAnnMin}{0.964}
\newcommand{\CAnnMax}{0.968}
\newcommand{\CDeltaMin}{0.496}
\newcommand{\CDeltaMax}{0.527}
\newcommand{\CRefDnbMin}{0.680}
\newcommand{\CRefDnbMax}{0.712}
\newcommand{\CRefHynMin}{0.556}
\newcommand{\CRefHynMax}{0.617}
\newcommand{\RevAnnMin}{0.965}
\newcommand{\RevAnnMax}{0.968}
\newcommand{\DnbSens}{0.174}
\newcommand{\NCExonCom}{8,339}
\newcommand{\NCExonPath}{0}
\newcommand{\DownstreamCom}{7,779}
\newcommand{\DownstreamPath}{0}
\newcommand{\UpstreamCom}{6,959}
\newcommand{\UpstreamPath}{0}
\newcommand{\RegulatoryCom}{5,247}
\newcommand{\RegulatoryPath}{0}
\newcommand{\NCHyphenCom}{0}
\newcommand{\NCHyphenPath}{2,996}
\newcommand{\InitTypoCom}{0}
\newcommand{\InitTypoPath}{663}
\newcommand{\RelHyn}{0.525}
\newcommand{\RelDnb}{4.05}
\newcommand{\RelKm}{6.33}
\newcommand{\RelAg}{0.00867}

\newcommand{\tbd}[1]{\textcolor{red}{\textbf{[#1]}}}
\newcommand{\Dlt}{\Delta}

\title{\textbf{Supplementary Material}\\[4pt]\large When Accuracy Is Not Understanding: Auditing Annotation Shortcuts in a ClinVar-gnomAD Benchmark for AI-Assisted Variant Interpretation}
\date{}
\begin{document}
\maketitle

\section*{Contents}
Table~S1 fold-level AUROC; Table~S2 probe regularisation; Table~S3 ClinVar review-status sensitivity; Table~S4 all annotation strata; Table~S5 incremental value of the reference-only score; Table~S6 consequence-term counts; Section~S7 reproducibility details and analyses not included.

\begin{table}[h]
\centering\small
\begin{tabular}{lccc}\toprule
Feature source & Group 1 & Group 2 & Group 3\\\midrule
Annotation: consequence terms + substitution & 0.969 & 0.970 & 0.965\\
Annotation: harmonised flags (2) & 0.873 & 0.828 & 0.858\\
Annotation: plain-missense flag only & 0.854 & 0.820 & 0.838\\
DNABERT-2, reference-only & 0.682 & 0.691 & 0.685\\
DNABERT-2, variant ($\Delta$) & 0.518 & 0.511 & 0.521\\
HyenaDNA, reference-only & 0.620 & 0.619 & 0.647\\
HyenaDNA, variant ($\Delta$) & 0.546 & 0.526 & 0.550\\
4-mer, reference-only & 0.566 & 0.579 & 0.582\\
4-mer, variant ($\Delta$) & 0.507 & 0.534 & 0.512\\
Substitution type only (no annotation) & 0.586 & 0.613 & 0.600\\
\bottomrule\end{tabular}

\caption{AUROC within each held-out chromosome group (pooled out-of-fold predictions split by group). Groups follow the definitions of Feng et al.\ (2025): group 1 = chr3, 6, 9, 12, 16, 18, 19, 21; group 2 = chr2, 5, 11, 14, 17, 20, 22, X; group 3 = chr1, 4, 7, 8, 10, 13, 15.}
\end{table}

\begin{table}[h]
\centering\small
\begin{tabular}{lrcc}\toprule
Feature source & $C$ & Reference-only AUROC & Variant ($\Delta$) AUROC\\\midrule
DNABERT-2 & 0.001 & 0.712 & 0.502\\
DNABERT-2 & 0.01 & 0.708 & 0.504\\
DNABERT-2 & 0.1 & 0.691 & 0.506\\
DNABERT-2 & 1 & 0.681 & 0.507\\
DNABERT-2 & 10 & 0.680 & 0.507\\
HyenaDNA & 0.001 & 0.556 & 0.511\\
HyenaDNA & 0.01 & 0.593 & 0.525\\
HyenaDNA & 0.1 & 0.612 & 0.527\\
HyenaDNA & 1 & 0.617 & 0.524\\
HyenaDNA & 10 & 0.616 & 0.523\\
4-mer & 0.001 & 0.573 & 0.496\\
4-mer & 0.01 & 0.574 & 0.504\\
4-mer & 0.1 & 0.572 & 0.504\\
4-mer & 1 & 0.568 & 0.504\\
4-mer & 10 & 0.564 & 0.504\\
Annotation (terms + substitution) & 0.001 & - & 0.964\\
Annotation (terms + substitution) & 0.01 & - & 0.968\\
Annotation (terms + substitution) & 0.1 & - & 0.968\\
Annotation (terms + substitution) & 1 & - & 0.968\\
Annotation (terms + substitution) & 10 & - & 0.967\\
\bottomrule\end{tabular}

\caption{AUROC across L2 regularisation strengths $C$ (inverse regularisation; larger is weaker). The main text uses $C=1$, fixed in advance and identical for all feature sources.}
\end{table}

\begin{table}[h]
\centering\small
\resizebox{\linewidth}{!}{\begin{tabular}{llcc}\toprule
Pathogenic subset (pathogenic / common) & Feature source & AUROC (95\% CI)\\\midrule
\multicolumn{3}{l}{\emph{All variants} (3,000 / 3,000)}\\
\quad Annotation: consequence terms + substitution & & 0.968 (0.963--0.972)\\
\quad DNABERT-2, reference-only & & 0.681 (0.668--0.694)\\
\quad DNABERT-2, variant ($\Delta$) & & 0.507 (0.492--0.522)\\
\quad HyenaDNA, reference-only & & 0.617 (0.604--0.632)\\
\multicolumn{3}{l}{\emph{Pathogenic with assertion criteria} (1,763 / 3,000)}\\
\quad Annotation: consequence terms + substitution & & 0.968 (0.963--0.972)\\
\quad DNABERT-2, reference-only & & 0.690 (0.675--0.705)\\
\quad DNABERT-2, variant ($\Delta$) & & 0.512 (0.495--0.530)\\
\quad HyenaDNA, reference-only & & 0.614 (0.598--0.629)\\
\multicolumn{3}{l}{\emph{Pathogenic multi-submitter/expert} (554 / 3,000)}\\
\quad Annotation: consequence terms + substitution & & 0.965 (0.959--0.969)\\
\quad DNABERT-2, reference-only & & 0.689 (0.667--0.711)\\
\quad DNABERT-2, variant ($\Delta$) & & 0.491 (0.464--0.518)\\
\quad HyenaDNA, reference-only & & 0.620 (0.594--0.644)\\
\bottomrule\end{tabular}
}
\caption{AUROC after restricting the pathogenic class by ClinVar review status (all common variants retained). The first subset is the full sample. Annotation-only AUROC is essentially unchanged, so the findings are not driven by the weakest assertions.}
\end{table}

\begin{table}[h]
\centering\footnotesize
\resizebox{\linewidth}{!}{\begin{tabular}{llcc}\toprule
Stratum (pathogenic / common) & Feature source & AUROC (95\% CI) & Probe trained\\\midrule
\multicolumn{4}{l}{\emph{Plain missense only} (2,247 / 236)}\\
\quad & HyenaDNA reference-only & 0.635 (0.596-0.674) & all variants\\
\quad & HyenaDNA variant ($\Delta$) & 0.527 (0.487-0.563) & all variants\\
\quad & DNABERT-2 reference-only & 0.749 (0.711-0.782) & all variants\\
\quad & DNABERT-2 variant ($\Delta$) & 0.493 (0.455-0.532) & all variants\\
\quad & 4-mer reference-only & 0.552 (0.512-0.591) & all variants\\
\quad & 4-mer variant ($\Delta$) & 0.460 (0.422-0.501) & all variants\\
\quad & Annotation & 0.582 (0.545-0.620) & all variants\\
\quad & HyenaDNA reference-only & 0.649 (0.609-0.688) & within stratum\\
\quad & DNABERT-2 reference-only & 0.721 (0.683-0.757) & within stratum\\
\midrule
\multicolumn{4}{l}{\emph{Non-coding-transcript flagged} (424 / 1,897)}\\
\quad & HyenaDNA reference-only & 0.585 (0.556-0.614) & all variants\\
\quad & HyenaDNA variant ($\Delta$) & 0.509 (0.478-0.540) & all variants\\
\quad & DNABERT-2 reference-only & 0.633 (0.607-0.662) & all variants\\
\quad & DNABERT-2 variant ($\Delta$) & 0.502 (0.472-0.532) & all variants\\
\quad & 4-mer reference-only & 0.553 (0.523-0.583) & all variants\\
\quad & 4-mer variant ($\Delta$) & 0.512 (0.483-0.542) & all variants\\
\quad & Annotation & 1.000 (1.000-1.000) & all variants\\
\quad & HyenaDNA reference-only & 0.561 (0.534-0.591) & within stratum\\
\quad & DNABERT-2 reference-only & 0.560 (0.528-0.587) & within stratum\\
\midrule
\multicolumn{4}{l}{\emph{Other annotated} (329 / 867)}\\
\quad & HyenaDNA reference-only & 0.604 (0.569-0.640) & all variants\\
\quad & HyenaDNA variant ($\Delta$) & 0.559 (0.522-0.594) & all variants\\
\quad & DNABERT-2 reference-only & 0.611 (0.577-0.644) & all variants\\
\quad & DNABERT-2 variant ($\Delta$) & 0.565 (0.527-0.601) & all variants\\
\quad & 4-mer reference-only & 0.549 (0.514-0.586) & all variants\\
\quad & 4-mer variant ($\Delta$) & 0.563 (0.523-0.604) & all variants\\
\quad & Annotation & 0.976 (0.969-0.984) & all variants\\
\quad & HyenaDNA reference-only & 0.574 (0.539-0.610) & within stratum\\
\quad & DNABERT-2 reference-only & 0.566 (0.531-0.605) & within stratum\\
\midrule
\bottomrule\end{tabular}
}
\caption{All annotation strata. ``Annotation'' within the non-coding-transcript stratum reaches AUROC 1.000 because the two classes spell the non-coding term differently; this is a vocabulary artefact, not a biological signal. Strata are defined from the consequence string: plain missense; any non-coding-transcript term; other.}
\end{table}

\begin{table}[h]
\centering\small
\begin{tabular}{lccc}\toprule
Model & Annotation alone & + reference-only score & Gain (95\% CI)\\\midrule
HyenaDNA & 0.968 & 0.972 & 0.005 (0.002-0.007)\\
DNABERT-2 & 0.968 & 0.977 & 0.010 (0.007-0.013)\\
4-mer & 0.968 & 0.968 & 0.001 (-0.001-0.002)\\
\bottomrule\end{tabular}

\caption{Does the reference-only embedding score add information beyond annotation (consequence terms plus substitution type)? Gain is the paired-bootstrap difference in AUROC between the combined model and annotation alone. The gains are statistically distinguishable from zero for DNABERT-2 and HyenaDNA but small.}
\end{table}

\begin{table}[h]
\centering\small
\begin{tabular}{lrrr}\toprule
Consequence term & Pathogenic (n) & Common (n) & Total\\\midrule
\texttt{missense\_variant} & 22,254 & 16,557 & 38,811\\
\texttt{non\_coding\_transcript\_exon\_variant} & 0 & 8,339 & 8,339\\
\texttt{downstream\_gene\_variant} & 0 & 7,779 & 7,779\\
\texttt{intron\_variant} & 1,702 & 5,325 & 7,027\\
\texttt{upstream\_gene\_variant} & 0 & 6,959 & 6,959\\
\texttt{regulatory\_region\_variant} & 0 & 5,247 & 5,247\\
\texttt{intron\_variant\&non\_coding\_transcript\_variant} & 0 & 4,736 & 4,736\\
\texttt{non-coding\_transcript\_variant} & 2,996 & 0 & 2,996\\
\texttt{5\_prime\_UTR\_variant} & 930 & 1,571 & 2,501\\
\texttt{3\_prime\_UTR\_variant\&NMD\_transcript\_variant} & 0 & 2,289 & 2,289\\
\texttt{missense\_variant\&NMD\_transcript\_variant} & 0 & 2,220 & 2,220\\
\texttt{intron\_variant\&NMD\_transcript\_variant} & 0 & 1,916 & 1,916\\
\texttt{3\_prime\_UTR\_variant} & 271 & 1,012 & 1,283\\
\texttt{synonymous\_variant} & 64 & 723 & 787\\
\texttt{TF\_binding\_site\_variant} & 0 & 732 & 732\\
\texttt{initiatior\_codon\_variant} & 663 & 0 & 663\\
\texttt{missense\_variant\&splice\_region\_variant} & 0 & 383 & 383\\
\texttt{splice\_region\_variant\&non\_coding\_transcript\_exon\_variant} & 0 & 220 & 220\\
\texttt{synonymous\_variant\&NMD\_transcript\_variant} & 0 & 181 & 181\\
\texttt{nonsense} & 104 & 0 & 104\\
\bottomrule\end{tabular}

\caption{The twenty most frequent consequence terms (counts of variants carrying the term) in the full annotation table. A variant can carry several terms.}
\end{table}

\section*{S7. Reproducibility details and analyses not included}
\textbf{Software and settings.} Python with numpy, pandas, SciPy, scikit-learn (logistic regression, L-BFGS; StandardScaler fitted on training folds) and matplotlib; random seed 42 for sampling, resampling and probes. Versions: Python 3.13, NumPy 2.5.3, pandas 3.0.5, SciPy 1.18.1, scikit-learn 1.9.1, Matplotlib 3.11.2; \texttt{alphagenome-pytorch} 0.3.1 (embeddings computed earlier).

\textbf{Files.} The analysis folder contains \texttt{scripts/01\_leakage\_main.py} (annotation-only and embedding probes, stratified analysis, incremental value, label-shuffle control), \texttt{02\_crosstab\_and\_alphagenome.py} (vocabulary cross-tabulation; AlphaGenome subset), \texttt{03\_shared\_vocab\_and\_figure.py} (shared-vocabulary probe; Figure~2), \texttt{04\_make\_manuscript\_assets.py} and \texttt{05\_extra\_assets.py} (robustness analyses and every table and number in the manuscript), and \texttt{build\_manuscript.sh}. Result tables are in \texttt{results/}. Per-variant out-of-fold scores are in \texttt{results/oof\_scores\_main.csv}.

\textbf{Data sources.} Variants and annotation: public benchmark of Feng et al.\ (2025) and its annotation table. Embeddings: DNABERT-2, HyenaDNA and a 4-mer baseline (501-bp window, \NPerClass{} variants per class, seed 42), and AlphaGenome embeddings for the \AgN{}-variant upstream test split. All embeddings were computed before this analysis.

\textbf{Counts.} The annotation table contains \TabNPath{} pathogenic and \TabNCommon{} unique common variants after removing duplicate coordinates; this differs by a few rows from the \(39{,}637\) variants in the source sequence table. The sampled \NSample{} variants were all matched to annotation.

\textbf{Analyses not included.} An earlier version of the project compared seven models with pairwise significance tests and an embedding-geometry analysis (principal components, silhouette, $k$-nearest-neighbour purity, centred kernel alignment, UMAP). These were removed because the study's question concerns benchmark construction, not model ranking, and because variant-specific results were uniformly close to chance. The relevant models were not rerun.

\textbf{Not done.} Pooling at the variant position, intermediate-layer analyses, longer-window replication, gene-held-out evaluation and any ancestry-stratified analysis were not performed.

\end{document}
